\documentclass[12pt]{report}
\usepackage[T1]{fontenc}
\usepackage[brazilian]{babel}
\usepackage{csquotes}
\usepackage[
  perfil=projeto,
  impressao=frente-e-verso,
  citacoes=autor-data
]{abntex3}

\addbibresource{referencias-exemplo.bib}

\ABNTEXsetup{
  titulo={Preservação digital de documentos acadêmicos},
  subtitulo={projeto de pesquisa do Marco 8},
  autoria={Maria da Silva},
  instituicao={Universidade Exemplo},
  natureza={Projeto de pesquisa},
  objetivo={submetido ao Programa de Pós-Graduação em Documentação},
  area={Preservação digital},
  orientacao={Prof. Dr. João Souza},
  local={Fortaleza},
  data={2026}
}

\begin{document}

\ABNTEXprojetoexterna

% A capa é opcional neste perfil.
\ABNTEXcapa

\ABNTEXprojetopretextual
\ABNTEXfolhaderosto
\ABNTEXsumario

\ABNTEXprojetotextual

\ABNTEXprojetointroducao{%
  O projeto investiga a preservação de documentos acadêmicos digitais.
  O problema consiste em identificar práticas reprodutíveis para reduzir
  perdas de informação. Parte-se da hipótese de que formatos abertos e
  verificações automatizadas aumentam a durabilidade dos acervos. O objetivo
  é comparar estratégias existentes e justificar um fluxo institucional.
}

\ABNTEXprojetoreferencialteorico{%
  A análise considera princípios de normalização documental e preservação
  digital, mantendo separadas a apresentação do projeto e a descrição
  bibliográfica processada pelo \texttt{biblatex-abnt}
  \parencite{associacao2025}.
}

\ABNTEXprojetometodologia{%
  Será realizado um estudo comparativo com amostras controladas, registro dos
  formatos, cálculo de somas de verificação e reprodução dos resultados.
}

\ABNTEXprojetorecursos{%
  Serão utilizados um repositório versionado, armazenamento redundante,
  ferramentas livres de validação e horas de trabalho da equipe.
}

\ABNTEXprojetocronograma{%
  A revisão bibliográfica ocorrerá no primeiro trimestre; a coleta e os
  ensaios, no segundo; a análise e a redação, no terceiro.
}

\ABNTEXprojetopostextual
\ABNTEXbibliografia

\end{document}
